\documentclass[11pt]{amsart}
\usepackage[margin=1.1in]{geometry}
\usepackage{amsmath,amssymb,amsthm,booktabs}
\usepackage[colorlinks=true,linkcolor=blue,citecolor=blue,urlcolor=blue]{hyperref}
\newtheorem{theorem}{Theorem}[section]
\newtheorem{proposition}[theorem]{Proposition}
\newtheorem{lemma}[theorem]{Lemma}
\theoremstyle{remark}
\newtheorem{remark}[theorem]{Remark}
\newcommand{\Z}{\mathbb Z}\newcommand{\Q}{\mathbb Q}\newcommand{\PP}{\mathbb P}\newcommand{\F}{\mathbb F}
\newcommand{\NS}{\operatorname{NS}}\newcommand{\MW}{\operatorname{MW}}
\title{An Ap\'ery-like sequence attached to a singular $K3$ surface}
\author{Claude (Fable), with River}
\date{23 August 2026 --- draft}
\begin{document}
\begin{abstract}
The integer sequence $1,4,16,64,250,928,3136,8704,11866,-79400,\dots$ defined by the four-term recurrence $(n+1)^2u_{n+1}=(11n^2+11n+4)u_n-(37n^2+3)u_{n-1}+3(3n-1)(3n-2)u_{n-2}$ is the period sequence of an elliptic $K3$ surface with five singular fibres $I_4,I_6,I_3,I_3,IV^*$ and Picard number $20$. Its transcendental motive is that of the weight-$3$ CM newform \texttt{32.3.d.a} (CM by $\Q(\sqrt{-2})$), and its Ap\'ery constants are critical $L$-values of that form: the fold constants at the conjugate singularities $\tfrac{5\mp i\sqrt2}{27}$ are $\tfrac{\sqrt2\mp i}3L(g,2)$ and the constant at $t=1$ is $\tfrac{2\sqrt2}3L(g,2)$, with $L(g,2)=\tfrac\pi{32}\Gamma(\tfrac18)\Gamma(\tfrac38)/\Gamma(\tfrac58)\Gamma(\tfrac78)$ by Chowla--Selberg. To our knowledge this is the first Ap\'ery-like sequence whose Ap\'ery constant is a critical value of a cusp form; three-term integral recurrences can never produce one, by a Riemann--Hurwitz/twist-parity bound that confines them to rational elliptic surfaces.
\end{abstract}
\maketitle

\section{The sequence and the results}
Let $u_{-1}=u_{-2}=0$, $u_0=1$ and
\begin{equation}\label{rec}
(n+1)^2u_{n+1}=(11n^2+11n+4)\,u_n-(37n^2+3)\,u_{n-1}+3(3n-1)(3n-2)\,u_{n-2}.
\end{equation}
All $u_n$ are integers (verified for $n\le300$; no proof is known --- see Remark~\ref{integrality}). The sequence arose in a systematic scan \cite{FTS} of four-term Ap\'ery-like recurrences: among $2178$ integral rows found in the Kodaira-admissible classes, $98\%$ are three-term rows in disguise (an apparent or $I_0^*$ singular point) and, of the $30$ genuinely five-point sporadic rows, \eqref{rec} is the unique one whose local system is the Picard--Fuchs system of an elliptic surface.

Write $A(t)=\sum_{n\ge0}u_nt^n$; it satisfies $LA=0$ for the self-adjoint operator
\[L\,y=\bigl(t\,R(t)\,y'\bigr)'-4(1-10t+15t^2)\,y,\qquad R(t)=1-11t+37t^2-27t^3=(1-t)(1-10t+27t^2),\]
with singular points $t=0,\ \tfrac{5\pm i\sqrt2}{27},\ 1,\ \infty$, all local exponents $(0,0)$ at the finite points. Let $B(t)=\sum b_nt^n$ be the companion solution of $LB=1$ with $b_0=0,b_1=1$; then $d_n^2b_n\in\Z$ with $d_n=\operatorname{lcm}(1,\dots,n)$, sharp.

\begin{theorem}\label{main}
\begin{enumerate}
\item[(i)] The local system of $L$ is $R^1\pi_*\Q$ of an elliptic $K3$ surface $\pi:\mathcal X\to\PP^1_t$ with singular fibres $I_4\,(t=0)$, $I_6\,(t=1)$, $I_3,I_3\,(t=\tfrac{5\pm i\sqrt2}{27})$, $IV^*\,(t=\infty)$, given by the explicit Weierstrass model $y^2=x^3-27c_4x-54c_6$ of \S2. The functional invariant $\mathcal J$ is a Belyi map of degree $16$ attaining the Riemann--Hurwitz bound.
\item[(ii)] $\rho(\mathcal X)=20$ \textup{(}a singular $K3$\textup{)}, $\MW$ has rank $0$ and torsion exactly $\Z/3$, generated by a section over $\Q(\sqrt3)$ with $x_0=9(4t-1)^2$; the transcendental lattice is $T\cong\operatorname{diag}(6,12)$, $\det T=72$, with CM by $\Q(\sqrt{-2})$.
\item[(iii)] The transcendental motive is that of the newform $g=\texttt{32.3.d.a}=\texttt{8.3.d.a}\otimes\chi_{-4}\in S_3(\Gamma_0(32),\chi_{-8})$, $\texttt{8.3.d.a}=\eta(z)^2\eta(2z)\eta(4z)\eta(8z)^2$: the traces $a_p=\sum_{t,x\in\F_p}\chi_p\bigl(x^3-27c_4(t)x-54c_6(t)\bigr)$ agree with $a_p(g)$ for all $76$ primes $5\le p\le397$ \textup{(}Livn\'e rigidity then identifies the motive\textup{)}.
\item[(iv)] The fold constants of the pair $(A,B)$ \textup{(}defined in \S3; verified to $165$--$230$ digits\textup{)} are
\[\xi\Bigl(\tfrac{5\mp i\sqrt2}{27}\Bigr)=\frac{\sqrt2\mp i}{3}\,L(g,2),\qquad
\xi(1)=\frac{2\sqrt2}{3}\,L(g,2)=\xi_++\xi_-,\qquad
L(g,2)=\frac{\pi}{32}\,\frac{\Gamma(\tfrac18)\Gamma(\tfrac38)}{\Gamma(\tfrac58)\Gamma(\tfrac78)} .\]
\item[(v)] $2$-adically $b_n/a_n\to0$ at logarithmic rate $v_2(b_n/a_n)=\log_2n+O(1)$; there is no $3$-adic limit.
\end{enumerate}
\end{theorem}
Items (i), (ii) are proved; (iii) is proved modulo the standard application of Livn\'e's theorem, with the finite verification stated; (iv) is a numerical identification at the stated precision, cross-validated by two independent methods; (v) is measured. Full details, scripts and logs: \cite{K3P,FTS}.

\begin{remark}
Every previously known Ap\'ery-like sequence with identified Ap\'ery constant computes a period of an \emph{Eisenstein} class --- $\zeta(2),\zeta(3)$, Dirichlet $L(2,\chi)$, $L(3,\chi)$ --- or, for the square-root rows, an $L$-value of a weight-$3$ CM form arising on a \emph{rational} elliptic surface host \cite{Beu87,Zag}. Theorem~\ref{main} gives the first Ap\'ery constant that is a critical cusp-form $L$-value attached to a $K3$. This is forced by the shape of the recurrence: for three-term rows the functional invariant satisfies $\deg\mathcal J\le12$ and a twist-parity argument confines the untwisted surface to Beauville's rational list \cite{CLASS}; five singular points allow $\deg\mathcal J\le18$, and \eqref{rec} realises a $K3$ with $\deg\mathcal J=16$.
\end{remark}

\section{The surface}
The scan's $\mathcal J$-map certification gives $\mathcal J=U/V$ with
\[V=t^4(t-1)^6(27t^2-10t+1)^3,\qquad U=-(4t-1)^3\bigl(216t^4+64t^3-48t^2+12t-1\bigr)^3,\]
\[U-1728V=-\bigl(5832t^8+34560t^7-30016t^6+12624t^5-4380t^4+1280t^3-240t^2+24t-1\bigr)^2,\]
fitted from $38$ coefficients of the canonical nome and verified on $73$ more ($h=4$, $\gamma=-1$: $\mathcal J(t)=j(-q(t)^4)$). Branch reconstruction (with the quadratic twist pinned by the exact series identities $c_4=E_4(Q)/y_0^4$, $c_6=E_6(Q)/y_0^6$ at $Q=-q(t)^4$) yields
\[c_4=(4t-1)\bigl(216t^4+64t^3-48t^2+12t-1\bigr),\qquad c_6=-\bigl(5832t^8+\cdots-1\bigr),\qquad \Delta=-V,\]
and Tate's algorithm reads off the fibre table of Theorem~\ref{main}(i); $\sum e=24$, so $\mathcal X$ is $K3$. Shioda--Tate gives $\rho\ge2+3+5+2+2+6=20=h^{1,1}$, hence $\rho=20$ and $\MW$ rank $0$. The $3$-division polynomial factors as $\psi_3=3(x-9(4t-1)^2)\cdot(\text{irreducible cubic})$ with $y_0^2=432(t-1)^4(27t^2-10t+1)^2$: a $3$-torsion section over $\Q(\sqrt3)$. The determinant count $\det T=648/m^2$ with the even-lattice condition $\det T\equiv0,3\pmod4$ and Shioda's injection $\MW_{\rm tors}\hookrightarrow\prod_vG_v$ pin $m=3$, $\det T=72$; the discriminant-form computation in \cite{K3P} identifies $T\cong\operatorname{diag}(6,12)=3\langle x^2+2y^2\rangle$, whence the CM field $\Q(\sqrt{-2})$.

For the trace identity in Theorem~\ref{main}(iii): with $\mathcal F=R^1\pi_*\Q_\ell$ on the good locus, $H^0_c=H^2_c=0$, and in $0\to\oplus_v\mathcal F^{I_v}\to H^1_c\to H^1(\PP^1,j_*\mathcal F)\to0$ the skyscraper contribution $\sum a_v$ cancels against the same terms in $-\sum_{t\in\PP^1(\F_p)}a_t$, leaving $\operatorname{tr}(\mathrm{Frob}_p\mid T\otimes\Q_\ell)$ on the nose (here $H^1(\PP^1,j_*\mathcal F)=T\otimes\Q_\ell$ because $\rho=20$ and the N\'eron--Severi classes are accounted). The bad points remain distinct mod every $p\ge5$ (the relevant resultants are supported on $\{2,3\}$; $\operatorname{disc}(27t^2-10t+1)=-8$). The match with \texttt{32.3.d.a} holds at all $76$ primes $5\le p\le397$ with no free parameter, and $a_p=0$ exactly at $p\equiv5,7\bmod8$, as CM by $\Q(\sqrt{-2})$ requires.

\section{The periods}
Each finite singular point has exponents $(0,0)$: a logarithmic fold. At such a point $t_c$ write
\[A=c_0w_0+c_1\bigl(w_0\log(t-t_c)+h\bigr),\qquad B=p+d_0w_0+d_1\bigl(w_0\log(t-t_c)+h\bigr),\qquad \xi(t_c):=d_1/c_1,\]
with $w_0$ the analytic Frobenius solution and $p$ the analytic particular solution; $\xi(t_c)$ is independent of branch and normalisation. When $t_c$ is the unique nearest singularity, $\xi(t_c)=\lim b_n/a_n$; here the nearest singularities are the conjugate pair $\tfrac{5\mp i\sqrt2}{27}$ ($|t|=27^{-1/2}$), so $b_n/a_n$ oscillates and the $\xi$'s are complex fold constants. They were computed two independent ways --- adaptive-step analytic continuation of the ODE with Frobenius matching at $230$ digits, and directly from the recurrence asymptotics --- and identified by integer-relation search against a basis containing $L(g,2)$ and its twists, $\pi^2\sqrt D$, and the discriminant-$(-8)$ Chowla--Selberg products; the identifications of Theorem~\ref{main}(iv) hold with residual below $10^{-165}$, and $L(g,2)$ itself matches the $\Gamma$-quotient at $230$ digits (Chowla--Selberg, as $g$ has CM by $\Q(\sqrt{-2})$; cf.\ \cite{ChS}). The relation $\xi(1)=\xi_++\xi_-$ is exact at the same precision. As a by-product the logarithm coefficients of $A$ itself are algebraic multiples of $1/\pi$: $c_1(1)=-2/(\pi\sqrt3)$, $c_1(t_+)=\tfrac1\pi(-\tfrac7{2\sqrt3}+\tfrac{4i}{\sqrt6})$.

\begin{remark}\label{integrality}
The integrality of $u_n$ is, at present, purely experimental ($n\le300$; also mod small primes to $n=10^4$ in the scan). For three-term rows integrality is explained by modular parametrisation (the Hauptmodul is a congruence-group uniformiser); here the monodromy group is an index-$16$ genus-$0$ subgroup attaining the Riemann--Hurwitz bound (signature data in \cite{K3P}), and we do not know whether it is congruence, nor a modular proof of integrality. This seems to us the most interesting open question the sequence raises: it is precisely Zagier's integrality phenomenon \cite{Zag} one step beyond the rigid (three-term) case, where the classification of \cite{CLASS} no longer forces modularity.
\end{remark}

\begin{remark}
$2$-adically $b_n/a_n\to0$: the cuspidal clause of the Euler-factor criterion of \cite{proj} ($\xi_p=0$ for cuspidal classes), observed here for the first time on a cusp-form row. The rate is logarithmic, not linear: the slope law for Eisenstein rows (Proposition~C of \cite{proj}) does not extend to cuspidal rows.
\end{remark}

\subsection*{Reproducibility}
All computations are in the repository \texttt{math-modular-sources}: the scan and the row's discovery in \texttt{consolidation/FOUR\_TERM\_SCAN.md} (\texttt{lattice/four\_term/}), the surface, newform and period work in \texttt{consolidation/K3\_ROW\_PERIOD.md} (\texttt{lattice/k3\_period/}), with PARI/GP and mpmath scripts and logs. The newform is \texttt{32.3.d.a} in the LMFDB.

\begin{thebibliography}{9}
\bibitem{Beu87} F.~Beukers, \emph{Irrationality proofs using modular forms}, Ast\'erisque 147--148 (1987), 271--283.
\bibitem{ChS} S.~Chowla, A.~Selberg, \emph{On Epstein's zeta-function}, J. reine angew. Math. 227 (1967), 86--110.
\bibitem{CLASS} \emph{Integral second-order Ap\'ery-like recurrences and rational elliptic surfaces}, companion paper (\texttt{paper/classification}), 2026.
\bibitem{FTS} \emph{Four-term Ap\'ery-like rows: the five-singular-point scan}, \texttt{consolidation/FOUR\_TERM\_SCAN.md}, 2026.
\bibitem{K3P} \emph{The period of the elliptic-$K3$ four-term row}, \texttt{consolidation/K3\_ROW\_PERIOD.md}, 2026.
\bibitem{Livne} R.~Livn\'e, \emph{Motivic orthogonal two-dimensional representations of $\mathrm{Gal}(\overline{\Q}/\Q)$}, Israel J. Math. 92 (1995), 149--156.
\bibitem{proj} \emph{Modular sources of Ap\'ery-like sequences} (main project paper), \texttt{consolidation/paper\_draft8.pdf}, 2026.
\bibitem{SI} T.~Shioda, H.~Inose, \emph{On singular $K3$ surfaces}, in: Complex Analysis and Algebraic Geometry, Cambridge (1977), 119--136.
\bibitem{Zag} D.~Zagier, \emph{Integral solutions of Ap\'ery-like recurrence equations}, in: Groups and Symmetries, CRM Proc. 47 (2009), 349--366.
\end{thebibliography}
\end{document}
